Let~$Q$ be an acyclic quiver with vertex set~$I$ and recall the symmetrizing matrix $D=\operatorname{diag}(\mathbf{d})$,
where $\mathbf{d}=(d_i: i\in I)$.
Write $n_{ij}=n_{ji}$ for the number of arrows connecting vertices $i,j\in I$ in~$Q$ and note that all such arrows point in the same direction.
We assume further that there are only f\/initely many arrows whose source or target is the vertex $i\in I$.
From the pair $(Q,\mathbf{d})$, called a~\emph{valued quiver}, we def\/ine an adjacency matrix $B=B_Q=(b_{ij})$ by the rule
\begin{gather*}
b_{ij}=
\begin{cases}
n_{ij}d_j/\gcd(d_i,d_j) & \text{if $i\to j$ in~$Q$,}
\\
-n_{ij}d_j/\gcd(d_i,d_j) & \text{if $j\to i$ in~$Q$,}
\\
0 & \text{if $i=j$.}
\end{cases}
\end{gather*}
Notice that the matrix~$B$ is skew-symmetrizable with skew-symmetrizing matrix~$D$, i.e.~$d_ib_{ij}=-d_jb_{ji}$ for all $i,j\in I$.
To connect with the previous section we will assume that the matrix~$B$ is related to the Cartan matrix~$A$~by
$a_{ij}=-|b_{ij}|$ for $i\ne j$.
